\documentclass[11pt]{article}

\usepackage[utf8]{inputenc}
\usepackage[T1]{fontenc}
\usepackage[english]{babel}
\usepackage{amsmath,amssymb,amsthm,mathtools}
\usepackage{booktabs}
\usepackage{microtype}
\usepackage{geometry}
\usepackage{enumitem}
\usepackage{xcolor}
\usepackage{hyperref}
\geometry{margin=1in}
\hypersetup{colorlinks=true,linkcolor=blue!55!black,citecolor=blue!55!black,urlcolor=blue!55!black}

\newtheorem{proposition}{Proposition}
\newtheorem{lemma}{Lemma}
\newtheorem{corollary}{Corollary}
\newtheorem{assumption}{Assumption}
\theoremstyle{definition}
\newtheorem{definition}{Definition}
\theoremstyle{remark}
\newtheorem{remark}{Remark}

\newcommand{\R}{\mathbb{R}}
\newcommand{\E}{\mathbb{E}}
\newcommand{\kb}{k_B}
\newcommand{\km}{k_M}
\newcommand{\Rb}{R_B}
\newcommand{\Rm}{R_M}

\title{Credit Quantity, Credit Prices, and Future Output:\\
A Model of Heterogeneous Financing Access}
\author{Alvaro Marcelo Chavez Unyen\\
\small Artificial Intelligence and Economic Modeling, Universidad del Pacifico}
\date{October 2026 -- research draft}

\begin{document}
\maketitle

\begin{abstract}
Preliminary evidence for Peru suggests that credit volume may contain more
forecasting information for non-tradable activity at relatively short
horizons, whereas a credit-price spread may become more informative for
tradable activity over a longer horizon. This paper develops a deliberately
small model that can rationalize that qualitative asymmetry. Firms invest at
date zero and produce at
date one. Bank credit is relatively cheap but capped; some firms can finance
investment beyond the bank ceiling in an external market at a wedge. The
firm's optimum has a quantity regime, in which the bank ceiling is binding,
and a price regime, in which external finance funds the marginal unit. In the
quantity regime, a larger credit ceiling raises investment and future output,
whereas a small spread change has no local real effect. In the price regime,
a larger bank ceiling changes financing composition but not investment, while
a spread increase lowers investment and output. Aggregating firms shows that
a sector with more bank-dependent firms is more quantity-sensitive and less
price-sensitive. The model provides a theoretical bridge from financing
heterogeneity to the sector-specific patterns being studied in the thesis. A
quadratic specialization yields closed-form first-order conditions. A direct
Lean 4 prototype checks the core algebra; the course-required
AppliedModelingLib formalization remains part of the final-paper workflow.
\end{abstract}

\noindent\textbf{Keywords:} credit constraints; external finance; sectoral
heterogeneity; investment; forecasting.\\
\textbf{JEL codes:} E22, E32, E44, G21, G32.

\section{Introduction}

Credit conditions can affect firms through two conceptually different
margins. A quantity restriction limits how much a firm can invest even when a
project has positive marginal value. A price change alters the marginal cost
of funds and therefore the desired scale of investment. Treating these margins
as interchangeable can hide economically meaningful heterogeneity.

Fernandez and Chavez (2026) compare autoregressive benchmarks with models
augmented by credit indicators for Peruvian tradable and non-tradable sectors.
Their work in progress points to a qualitative asymmetry: credit volume appears
more informative for non-tradable activity at relatively short horizons,
whereas a credit-price spread appears more informative for tradable activity
over a longer horizon. These results remain preliminary and should not be read
as a settled ranking of exact forecast horizons. Moreover, predictive content
alone does not supply an optimizing model that explains why the relevant
credit margin might differ across sectors.

This paper supplies that missing mechanism. Firms have internal funds and a
bank credit line. Bank credit is cheaper than external market finance, but the
line is capped. Firms differ in whether, and at what wedge, they can use an
alternative source after exhausting the bank line. Current financing choices
fund capital that produces next-period output, so the model generates a direct
lead-lag relation from credit conditions to real activity.

The main result is a regime separation. When the firm's desired capital stock
lies at the bank ceiling, quantity is locally decisive: relaxing the ceiling
raises capital one for one, while a small funding-price movement leaves capital
unchanged. When the firm uses market finance at the margin, price is decisive:
the common spread changes desired capital, while a larger bank line merely
replaces market finance with bank finance. Sectoral differences follow from
composition. If non-tradables contain more bank-dependent firms than tradables,
non-tradable output is more sensitive to credit quantity and tradable output is
more sensitive to spreads.

The contribution is intentionally narrow. It is not a full financial
accelerator or a structural interpretation of the thesis's forecast tests.
Instead, it isolates the minimum mechanism needed to convert the empirical
price-versus-quantity asymmetry into a proposition with explicit conditions.

\section{Related literature and positioning}

The first strand studies equilibrium credit rationing. Stiglitz and Weiss
(1981) show that asymmetric information can prevent the interest rate from
clearing the loan market, so observationally similar borrowers may face a
quantity restriction. Petersen and Rajan (1994) find that lending relationships
benefit small firms primarily through greater availability of funds rather
than large price reductions. The quantity regime below takes a credit ceiling
as a reduced-form representation of those frictions; it does not reproduce
their screening or relationship formation models.

The second strand studies balance-sheet and collateral amplification.
Kiyotaki and Moore (1997) make durable assets both productive inputs and
collateral, generating persistent responses through endogenous credit limits.
Bernanke, Gertler, and Gilchrist (1999) embed an external-finance premium and
heterogeneous capital-market access in a dynamic macroeconomic framework.
Holmstrom and Tirole (1997) link firm net worth and intermediary capital to the
supply of loanable funds. The present paper treats net worth and the credit
ceiling as predetermined in its baseline. Section~\ref{sec:extensions} shows
how an endogenous ceiling can recover an accelerator without obscuring the
main regime result.

The third strand establishes firm heterogeneity. Gertler and Gilchrist (1994)
document that small firms contract more sharply after monetary tightening,
while Petersen and Rajan (1994) emphasize their dependence on institutional
lenders. These findings motivate a higher constrained share among
non-tradable firms. This is a composition assumption to be disciplined by
data, not a theorem that every small firm is non-tradable or every tradable
firm has market access.

The fourth strand concerns leading financial indicators. Gilchrist and
Zakrajsek (2012) construct a credit-spread measure with predictive power for
future activity, and Faust et al. (2013) find that credit spreads improve
real-time activity forecasts. Stock and Watson (2003) caution that predictive
performance can be unstable across countries and periods. For Peru, Lahura
(2011) finds a stable long-run credit-output relationship, while Butron and
Pozo (2025) document a positive relationship using sectoral and regional
heterogeneity.

The closest paper for the sectoral motivation is Muller and Verner (2024).
They show that credit allocation toward non-tradable rather than tradable
sectors predicts different macroeconomic outcomes, discuss smaller and more
constrained non-tradable firms, and report robustness involving bond issuance.
That evidence means neither the tradable/non-tradable split nor financing
heterogeneity can be claimed as novel here. Holmstrom and Tirole's (1994)
working-paper version and their published 1997 model also already contain
capital thresholds associated with monitored and direct finance. Gilchrist,
Sim, and Zakrajsek (2014) subsequently model the joint dynamics of investment
and credit spreads.

The narrower course-paper contribution is to combine a capped, relatively
cheap bank tranche with a marginal outside-finance tranche and derive a
specific comparative-static ranking: a sector with more firms at the cap is
more sensitive to observed credit quantity, while a sector with more firms
financing the marginal unit externally is more sensitive to a common funding
price. The reviewed papers do not state that quantity-versus-price ranking in
this form. This is a provisional positioning claim, not a claim that the
broader ingredients or empirical sector pattern are new.

\section{Environment}

There are two dates, $t$ and $t+1$, and sectors $s\in\{N,T\}$, interpreted as
non-tradable and tradable. A firm enters date $t$ with internal funds $n>0$.
It chooses capital $k$, bank borrowing $b$, and external market finance $m$.
Capital becomes productive at $t+1$, making date-$t$ credit conditions leading
indicators of output.

\subsection{Technology}

The general model uses a differentiable production function $F:\R_+\to\R_+$.

\begin{assumption}[Technology]\label{ass:technology}
$F(0)=0$, $F'(k)>0$, and $F''(k)<0$ on the economically relevant domain.
Future revenue is $F(k)$ and its date-$t$ value is $\beta F(k)$, where
$\beta\in(0,1]$.
\end{assumption}

Closed forms and the Lean verification use
\begin{equation}\label{eq:quadratic}
F(k)=ak-\frac{c}{2}k^2,\qquad a>0,\quad c>0,
\end{equation}
restricted to $k<a/c$, where output is increasing. This quadratic form is a
local second-order representation of any smooth concave technology and keeps
all regime boundaries transparent.

\subsection{Financing}

The firm's sources satisfy
\begin{equation}\label{eq:financeidentity}
k=n+b+m,\qquad 0\le b\le \bar b,\qquad m\ge 0.
\end{equation}
Bank finance has gross marginal cost
\begin{equation}
\Rb(z)=R_f+z,
\end{equation}
where $z$ is a common credit-spread factor. External market finance costs
\begin{equation}
\Rm(z,\tau)=\Rb(z)+\tau,\qquad \tau\ge 0.
\end{equation}
The access wedge $\tau$ includes issuance costs, information costs, and other
barriers to non-bank funding. A bank-only firm is the limiting case in which
$m=0$ is imposed. A market-access firm can choose $m\ge0$.

\begin{assumption}[Funding order]\label{ass:fundingorder}
$0<\Rb\le\Rm<\beta a$. The firm uses internal funds first, bank credit second,
and market finance last. The strict upper bound guarantees a positive target
capital stock in the quadratic specialization.
\end{assumption}

The assumption that $z$ shifts both funding costs is important. Empirically,
the observed corporate loan spread is interpreted as a proxy for a broader
external-finance condition, not as an isolated price that leaves bond and
foreign-credit costs fixed. Section~\ref{sec:extensions} discusses this
identification boundary.

\section{The firm problem}

The firm maximizes discounted profits:
\begin{equation}\label{eq:firmproblem}
\max_{k,b,m}\quad \beta F(k)-\Rb b-\Rm m
\quad\text{s.t.}\quad \eqref{eq:financeidentity}.
\end{equation}
Let the capital stock attainable after exhausting the bank line be
\begin{equation}
\bar k\equiv n+\bar b.
\end{equation}
Because bank finance is weakly cheaper, the reduced problem is continuous and
piecewise:
\begin{equation}\label{eq:piecewiseprofit}
\Pi(k)=
\begin{cases}
\beta F(k)-\Rb(k-n), & n\le k\le\bar k,\\[3pt]
\beta F(k)-\Rb\bar b-\Rm(k-\bar k), & k>\bar k.
\end{cases}
\end{equation}
The left and right first-order conditions are
\begin{align}
\beta F'(k)&=\Rb &&\text{(bank-finance margin)},\label{eq:focbank}\\
\beta F'(k)&=\Rm &&\text{(market-finance margin)}.\label{eq:focmarket}
\end{align}
Let $\kb$ and $\km$ solve \eqref{eq:focbank} and
\eqref{eq:focmarket}. Strict concavity makes both targets unique. Since
$\Rm\ge\Rb$, declining marginal product implies $\km\le\kb$.

Under \eqref{eq:quadratic},
\begin{equation}\label{eq:targets}
\kb=\frac{\beta a-\Rb}{\beta c},
\qquad
\km=\frac{\beta a-\Rm}{\beta c},
\qquad \km\le\kb.
\end{equation}

\begin{lemma}[Optimal capital with market access]\label{lem:policy}
Under Assumptions~\ref{ass:technology}--\ref{ass:fundingorder}, the unique
optimal capital stock is
\begin{equation}\label{eq:policy}
k^*=\begin{cases}
n, & \kb\le n,\\
\kb, & n<\kb<\bar k,\\
\bar k, & \km\le\bar k\le\kb,\\
\km, & \bar k<\km.
\end{cases}
\end{equation}
For a bank-only firm, the last segment is infeasible and
$k^*=\min\{\max\{n,\kb\},\bar k\}$.
\end{lemma}

\begin{proof}
On $[n,\bar k]$, the derivative of profit is
$\beta F'(k)-\Rb$; on $(\bar k,\infty)$ it is
$\beta F'(k)-\Rm$. Both are strictly decreasing. If $\kb\le n$, the
left derivative is already non-positive at $n$. If
$n<\kb<\bar k$, it crosses zero at $\kb$ before the kink. If
$\km\le\bar k\le\kb$, the derivative is weakly positive immediately to the
left of the kink and weakly negative immediately to its right, so the kink is
optimal. Finally, if $\bar k<\km$, profit continues increasing after the bank
line is exhausted and reaches its maximum at $\km$. Strict concavity gives
uniqueness, with adjacent formulas agreeing at equality boundaries.
\end{proof}

\section{Quantity and price transmission}

The two empirically relevant cases are the kink and market-finance regimes.

\begin{proposition}[Quantity versus price]\label{prop:regimes}
Suppose all inequalities defining a regime are strict, so sufficiently small
parameter changes do not cause a regime switch.

\begin{enumerate}[label=(\roman*)]
\item If $\km<\bar k<\kb$, the bank ceiling binds and
\begin{equation}
k^*=\bar k=n+\bar b,\qquad
\frac{\partial k^*}{\partial\bar b}=1,\qquad
\frac{\partial k^*}{\partial z}=0.
\end{equation}
Hence
\begin{equation}\label{eq:quantityeffects}
\frac{\partial y_{t+1}}{\partial\bar b}=F'(\bar k)>0,
\qquad
\frac{\partial y_{t+1}}{\partial z}=0.
\end{equation}

\item If $\bar k<\km$, the firm uses market finance at the margin and
\begin{equation}
k^*=\km(z,\tau),\qquad
\frac{\partial k^*}{\partial\bar b}=0,\qquad
\frac{\partial k^*}{\partial z}=\frac{1}{\beta F''(\km)}<0.
\end{equation}
Consequently,
\begin{equation}\label{eq:priceeffects-general}
\frac{\partial y_{t+1}}{\partial\bar b}=0,
\qquad
\frac{\partial y_{t+1}}{\partial z}
=\frac{F'(\km)}{\beta F''(\km)}<0.
\end{equation}
A larger bank line replaces market finance one for one until $m=0$; it raises
profits by reducing funding costs but does not change capital or output.
\end{enumerate}
\end{proposition}

\begin{proof}
Part (i) follows from the third line of \eqref{eq:policy}. Since strict
inequalities preserve the regime locally, $k^*=n+\bar b$ is independent of
$z$. Differentiating $F(k^*)$ gives \eqref{eq:quantityeffects}.

For part (ii), the market-margin condition is
$\beta F'(\km)=\Rm(z,\tau)$. Implicit differentiation with respect to $z$,
using $\partial\Rm/\partial z=1$, yields
$\beta F''(\km)\,\partial\km/\partial z=1$. Because $F''<0$, the response is
negative. While $m>0$, a marginal increase in $\bar b$ increases bank
borrowing and reduces market finance by the same amount, leaving total capital
unchanged. Applying the chain rule to output gives
\eqref{eq:priceeffects-general}.
\end{proof}

For the quadratic technology, the price-regime comparative statics become
especially transparent:
\begin{equation}\label{eq:quadraticcs}
\frac{\partial \km}{\partial z}=-\frac{1}{\beta c},
\qquad
\frac{\partial F(\km)}{\partial z}
=-\frac{\Rm}{\beta^2c}<0.
\end{equation}

\begin{remark}[Local zero effects]
The zero derivatives are regime statements, not universal claims. A sufficiently
large price shock can move a constrained firm out of the kink, and a sufficiently
large increase in the bank ceiling can eliminate market borrowing. The model
therefore predicts threshold nonlinearities.
\end{remark}

\section{Sectoral aggregation and empirical mapping}

Let $\lambda_s\in[0,1]$ be the share of sector-$s$ activity generated by firms
in the quantity regime. For a transparent aggregation, suppose firms share
technology and financing parameters conditional on regime. Sectoral output is
\begin{equation}\label{eq:aggregate}
Y_{s,t+1}=\lambda_s F(n+\bar b_s)
+(1-\lambda_s)F\bigl(\km(z_s,\tau_s)\bigr).
\end{equation}

\begin{assumption}[Sector composition]\label{ass:composition}
The non-tradable sector contains a larger share of bank-dependent,
quantity-constrained firms:
\begin{equation}
0\le\lambda_T<\lambda_N\le1.
\end{equation}
\end{assumption}

\begin{proposition}[Sectoral sensitivity ranking]\label{prop:sector}
Under Assumption~\ref{ass:composition} and common conditional primitives,
\begin{align}
\frac{\partial Y_N}{\partial\bar b}
&>\frac{\partial Y_T}{\partial\bar b}\ge0,\label{eq:quantorder}\\
\left|\frac{\partial Y_N}{\partial z}\right|
&<\left|\frac{\partial Y_T}{\partial z}\right|.\label{eq:priceorder}
\end{align}
In the quadratic specialization,
\begin{align}
\frac{\partial Y_s}{\partial\bar b}
&=\lambda_s\bigl[a-c(n+\bar b)\bigr],\label{eq:aggquantity}\\
\left|\frac{\partial Y_s}{\partial z}\right|
&=(1-\lambda_s)\frac{\Rm}{\beta^2c}.
\label{eq:aggprice}
\end{align}
\end{proposition}

\begin{proof}
Differentiate \eqref{eq:aggregate} and apply
Proposition~\ref{prop:regimes}. The quantity response is the positive
within-regime marginal product multiplied by $\lambda_s$. The absolute price
response is the positive within-regime magnitude multiplied by
$1-\lambda_s$. Therefore $\lambda_N>\lambda_T$ establishes both strict
rankings. Equations \eqref{eq:aggquantity}--\eqref{eq:aggprice} substitute the
quadratic derivatives.
\end{proof}

This proposition offers a candidate interpretation for the thesis. If
bank-dependent small firms are more prevalent in non-tradable activities,
variation in available bank credit contains more information about their
feasible investment and subsequent production. If tradable firms are more
likely to access alternative sources, observed bank quantity can change
through substitution without a corresponding change in total finance. A broad
spread may then be more informative about their marginal funding condition.

The model delivers a structural sign restriction for a forecasting equation,
not a causal interpretation of an ADL coefficient:
\begin{equation}\label{eq:forecastmapping}
\Delta\log Y_{s,t+h}
\approx \eta^B_s\Delta\log\bar b_{s,t}
-\eta^z_s\Delta z_{s,t}+u_{s,t+h},
\end{equation}
where $\eta^B_N>\eta^B_T\ge0$ and
$\eta^z_T>\eta^z_N\ge0$. These inequalities are theoretical sign and ranking
restrictions, not estimates. The two-date model distinguishes a near-term
quantity channel from a potentially more persistent price channel but does not
identify exact forecast horizons. Doing so would require gestation lags,
adjustment costs, or a dynamic capital stock and is an extension rather than a
hidden assumption.

\section{Transmission mechanism and qualitative predictions}

The timing is central to the forecasting interpretation. At date $t$, the
bank ceiling and the common funding factor are observed before capital becomes
productive at $t+1$. They can therefore predict future activity even when the
model makes no causal claim about the reduced-form forecasting coefficient.
The predictive content comes from different sufficient statistics in the two
regimes.

\paragraph{Quantity-constrained firms.}
At the kink, the unconstrained bank-financed target lies above the cap, while
the market-financed target lies below it. The positive left derivative of
profits means that an additional unit of cheap bank funding would be installed
immediately. Observed bank credit is therefore informative about the feasible
investment pipeline, and the multiplier in \eqref{eq:quantityeffects} is the
marginal product of the capital unlocked by the line. A small spread movement
does not change capital while the firm remains at the kink; its first-order
real effect is locally muted.

\paragraph{Market-access firms.}
Once outside finance funds the marginal unit, bank lending is no longer a
sufficient statistic for total finance. A change in the bank line can induce a
one-for-one substitution between bank and market funding with no change in
capital. The common funding factor instead enters the marginal first-order
condition directly. It changes the desired capital stock and, with installation
or project-completion lags, can affect output over a more extended horizon.
This mechanism requires the measured spread to proxy for the relevant common
external-finance condition; it does not require literal bond issuance by every
tradable firm.

\paragraph{Qualitative empirical implications.}
The model suggests three checks rather than a mechanical mapping from theory
to forecast peaks. First, quantity predictive content should be stronger where
bank dependence or the frequency of exhausted credit lines is higher. Second,
spread predictive content should be stronger where firms have more alternative
funding access and their marginal cost co-moves with the measured spread.
Third, an expansion of bank lines should raise total investment for firms at
the cap but mainly change financing composition for firms already using
outside finance. These cross-sectional interactions are sharper tests of the
mechanism than any single preferred forecast horizon.

\section{Numerical illustration}

Consider $\beta=0.95$, $a=2$, $c=0.5$, $n=0.4$,
$\bar b=0.5$, $\Rb=0.5$, and $\Rm=1.4$. Then
\begin{equation}
\bar k=0.9,\qquad \kb=2.947,\qquad \km=1.053.
\end{equation}
A bank-only firm is quantity constrained at $k=0.9$. A market-access firm
uses $0.5$ of bank credit and approximately $0.153$ of market finance to reach
$k=1.053$.

\begin{table}[ht]
\centering
\caption{Illustrative regime responses}
\begin{tabular}{lcc}
\toprule
Experiment & Quantity-regime firm & Price-regime firm\\
\midrule
Baseline capital & 0.900 & 1.053\\
Baseline output & 1.598 & 1.828\\
$\bar b$ rises by 0.10 & capital rises to 1.000 & capital unchanged\\
$z$ rises by 0.02 & capital unchanged & capital falls to 1.011\\
Financing adjustment & more total bank credit & bank/market substitution\\
\bottomrule
\end{tabular}
\end{table}

With $\lambda_N=0.8$ and $\lambda_T=0.2$, the marginal quantity responses are
$1.24$ and $0.31$, respectively. The absolute spread responses are about
$0.62$ and $2.48$. The values are illustrative comparative statics, not a
calibration to Peru.

\section{Extensions and limitations}\label{sec:extensions}

\paragraph{Endogenous bank ceiling.}
Let $\bar b_{i,t}=\phi_s n_{i,t}$, with pledgeability $\phi_s>0$. In the
quantity regime, $k_{i,t}=(1+\phi_s)n_{i,t}$, so a net-worth shock changes
investment directly and through borrowing capacity. Adding
$n_{i,t+1}=\omega\Pi_{i,t}$ creates a simple financial accelerator. This
extension is useful for propagation, but not necessary for the core
price-versus-quantity proposition.

\paragraph{Continuous market access.}
The baseline binary distinction can be replaced by heterogeneous wedges
$\tau_i$ or finite outside-finance limits $\bar m_i$. Lower wedges make
$\bar k<\km$ more likely; higher wedges make the kink regime more likely.
The sectoral share $\lambda_s$ can therefore be derived from a distribution of
net worth, pledgeability, and access wedges.

\paragraph{Dynamic capital and horizons.}
A disciplined dynamic version introduces
$K_{t+1}=(1-\delta)K_t+I_t$, a convex installation cost $\Phi_s(I_t)$, and
the same two financing tranches for investment.  Let $q_{i,t}$ denote the
current-value shadow price of installed capital.  The investment and capital
first-order conditions imply
\begin{align}
 q_{i,t}
   &= R_{j,t}+\Phi_s'(I_{i,t})+\mu_{i,t},
   \label{eq:q-investment}\\
 q_{i,t}
   &= \beta\E_t\!\left[
       F_s'(K_{i,t+1})+(1-\delta)q_{i,t+1}
     \right],
   \label{eq:q-capital}
\end{align}
where $j=B$ when the marginal unit is bank financed and $j=M$ when it is
market financed.  The term $\mu_{i,t}\geq0$ is the shadow wedge attached to a
binding quantity limit; it is zero away from that limit (with the bank/market
cost gap instead determining which tranche is marginal).  Combining
\eqref{eq:q-investment}--\eqref{eq:q-capital} yields the Euler equation
\begin{equation}
 R_{j,t}+\Phi_s'(I_{i,t})+\mu_{i,t}
 =\beta\E_t\!\left[
 F_s'(K_{i,t+1})+(1-\delta)
 \{R_{j',t+1}+\Phi_s'(I_{i,t+1})+\mu_{i,t+1}\}
 \right].
 \label{eq:dynamic-euler}
\end{equation}
Thus a credit-quantity innovation operates through the current and expected
shadow wedges, whereas a spread innovation changes the marginal financing
price directly. Sector-specific curvature in $\Phi_s$, depreciation, or
project gestation lags can shift the output response across horizons and may
rationalize why preliminary quantity evidence appears at relatively short
horizons while price information appears over a longer horizon. Estimating
those dynamic parameters is left outside the current proposition so that the
first deliverable remains exactly provable by hand; equation
\eqref{eq:dynamic-euler} states the disciplined route for the next version
rather than claiming to identify the observed horizons already.

\paragraph{Measurement.}
The model's $z$ is a common marginal-finance factor. The thesis uses a bank
loan spread assigned by firm size because sector-level rates are unavailable.
The empirical proxy may therefore contain measurement error and need not move
one for one with bond or foreign-credit costs. A stronger empirical design
would test whether the spread predicts substitution across funding sources.

\paragraph{Boundaries of the result.}
The model abstracts from default, risk choice, lender optimization, general
equilibrium, and feedback from future expected output to current credit supply.
It explains a conditional transmission mechanism, not the origin of the
credit ceiling or the causal identification of credit shocks.

\section{Conclusion}

The same observed credit system can transmit shocks through different margins.
For a bank-dependent firm at a binding ceiling, additional credit raises
investment and future output. For a firm that can finance the marginal unit in
an external market, a larger bank line changes the source of funds but not the
scale of investment; the common funding spread governs that scale. A higher
share of the first type makes non-tradable activity quantity-sensitive, while
a higher share of the second makes tradable activity price-sensitive.

This mechanism turns the thesis's preliminary forecasting pattern into a
formal economic hypothesis with explicit assumptions, first-order conditions, regime
boundaries, and falsifiable predictions. Its value for the course is not that
it replaces the empirical thesis, but that it supplies the model the empirical
exercise previously treated as a black box.

\appendix
\section{Additional derivations}

For the quadratic technology, profit on a segment with marginal funding cost
$R$ is, up to constants,
\begin{equation}
\pi_R(k)=\beta\left(ak-\frac c2k^2\right)-Rk.
\end{equation}
The FOC and second-order condition are
\begin{equation}
\pi_R'(k)=\beta(a-ck)-R=0,\qquad
\pi_R''(k)=-\beta c<0.
\end{equation}
Completing the square gives the global identity
\begin{equation}
\pi_R(k_R)-\pi_R(k)=\frac{\beta c}{2}(k-k_R)^2\ge0,
\end{equation}
where $k_R=(\beta a-R)/(\beta c)$. At the bank ceiling, global optimality in
the kink regime follows from
\begin{align}
\Pi(\bar k)-\Pi(k)
&=(\bar k-k)\bigl[\beta F'(\bar k)-\Rb\bigr]
+\frac{\beta c}{2}(\bar k-k)^2\ge0,
&&k\le\bar k,\\
\Pi(\bar k)-\Pi(k)
&=(k-\bar k)\bigl[\Rm-\beta F'(\bar k)\bigr]
+\frac{\beta c}{2}(k-\bar k)^2\ge0,
&&k\ge\bar k.
\end{align}
These are the identities formalized by the Lean theorem
\texttt{kink\_regime\_optimal}.

\section{Cobb-Douglas robustness}

If $F(k)=Ak^\alpha$ with $A>0$ and $0<\alpha<1$, the two unconstrained targets
are
\begin{equation}
k_j=\left(\frac{\beta\alpha A}{R_j}\right)^{\frac{1}{1-\alpha}},
\qquad j\in\{B,M\}.
\end{equation}
The policy function retains exactly the four cases in
Lemma~\ref{lem:policy}. In the market-finance regime,
\begin{equation}
\frac{\partial k_M}{\partial z}
=-\frac{k_M}{(1-\alpha)R_M}<0,
\qquad
\frac{\partial F(k_M)}{\partial z}
=-\frac{\alpha}{1-\alpha}\frac{F(k_M)}{R_M}<0.
\end{equation}
Thus the regime ranking is not an artifact of the quadratic technology.

\section{Lean formalization status}

Table~\ref{tab:lean} records the declaration-level status of every numbered
result. The current files in \texttt{lean/}\allowbreak\texttt{prototype/} are a direct Lean 4
prototype: they verify the main quadratic identities and sensitivity
inequalities without \texttt{sorry}, but they are not yet the complete folder
generated by the course's Applied\-Modeling\-Lib workflow. Before the final paper,
the merged paper PDF will be pinned to a commit, the workflow will be run on
that version, and the entire generated folder will replace the prototype.

\begin{table}[ht]
\centering
\small
\caption{Paper result to Lean declaration map}
\label{tab:lean}
\begin{tabular}{@{}p{0.20\linewidth}p{0.30\linewidth}p{0.40\linewidth}@{}}
\toprule
Paper result & Prototype declaration(s) & Current status \\
\midrule
Lemma~\ref{lem:policy} & \texttt{kink\_regime\_}\newline
\texttt{optimal}; \texttt{market\_regime\_}\newline
\texttt{optimal} & Core kink and market branches proved;
full four-case wrapper pending Applied\-Modeling\-Lib run. \\
Proposition~\ref{prop:regimes} & \texttt{quantity\_output\_gain};
\texttt{market\_target\_}\newline
\texttt{decreases} & Proved for the quadratic
specialization and strict local regimes. \\
Proposition~\ref{prop:sector} & \texttt{quantity\_sensitivity\_}\newline
\texttt{order}; \texttt{price\_sensitivity\_}\newline
\texttt{order} & Proved under $\lambda_T<\lambda_N$ and
the positivity assumptions stated in the paper. \\
\bottomrule
\end{tabular}
\end{table}

Lean verifies mathematics conditional on supplied assumptions. It does not
verify that those assumptions hold in Peruvian data, that the literature
review is exhaustive, or that the empirical estimates identify causal
effects. Paper commit formalized: \textbf{pending}; Applied\-Modeling\-Lib
\texttt{check --fast}: \textbf{pending}.

\section{AI collaboration log}

AI assistance was used to search and compare the literature, propose the
two-tranche financing model, check symbolic and numerical claims, draft LaTeX,
and prototype Lean statements. The author retained responsibility for the
economic question, the interpretation of the thesis evidence, and each
accepted assumption. Raw prompts and relevant answers are preserved in
\texttt{prompts.md}.

\begin{table}[ht]
\centering
\small
\caption{Claims checked during AI-assisted development}
\begin{tabular}{p{0.29\linewidth}p{0.27\linewidth}p{0.34\linewidth}}
\toprule
AI-assisted claim & Independent check & Verdict \\
\midrule
A binding cheap-bank cap makes capital locally quantity-sensitive and
price-insensitive. & Hand derivation, regime inequalities, and
\texttt{code/verify.py}. & Accepted only for strict local kink regimes; large
shocks can switch regimes. \\
Outside finance makes the spread relevant and the bank line irrelevant for
total capital. & FOC, implicit derivative, numerical check. & Accepted while
market finance remains strictly positive; bank credit still affects profits
and funding composition. \\
Sector composition generates the tradable/non-tradable ranking. & Direct
differentiation of \eqref{eq:aggregate} and literature audit. & Accepted as a
conditional prediction, not as proof that the composition assumption holds in
Peru. \\
The model explains exact empirical forecast horizons. & Compared the static
timing with the unfinished thesis results. & Rejected. The draft now states
only a qualitative short-versus-longer-horizon interpretation. \\
The sector pattern itself is novel. & Checked Muller and Verner (2024), its
appendices, and related working papers. & Rejected. Novelty is limited to the
explicit quantity-versus-price comparative-static ranking. \\
\bottomrule
\end{tabular}
\end{table}

\begin{thebibliography}{99}

\bibitem{BGG1999}
Bernanke, B. S., Gertler, M., and Gilchrist, S. (1999).
The financial accelerator in a quantitative business cycle framework.
In J. B. Taylor and M. Woodford (eds.), \emph{Handbook of Macroeconomics},
Vol. 1C, 1341--1393.

\bibitem{ButronPozo2025}
Butron, N., and Pozo, J. (2025).
Dynamic relationship between credit and economic growth.
BCRP Working Paper 2025-011.

\bibitem{Thesis2026}
Fernandez, S. A., and Chavez, A. M. (2026).
\emph{Credit Indicators and Sectoral Activity in Peru: Evidence on Predictive
Power, 2008--2024}. Undergraduate research thesis, Universidad del Pacifico.

\bibitem{FaustEtAl2013}
Faust, J., Gilchrist, S., Wright, J. H., and Zakrajsek, E. (2013).
Credit spreads as predictors of real-time economic activity: A Bayesian
model-averaging approach. \emph{Review of Economics and Statistics},
95(5), 1501--1519.

\bibitem{GertlerGilchrist1994}
Gertler, M., and Gilchrist, S. (1994).
Monetary policy, business cycles, and the behavior of small manufacturing
firms. \emph{Quarterly Journal of Economics}, 109(2), 309--340.

\bibitem{GilchristZakrajsek2012}
Gilchrist, S., and Zakrajsek, E. (2012).
Credit spreads and business cycle fluctuations.
\emph{American Economic Review}, 102(4), 1692--1720.

\bibitem{GilchristSimZakrajsek2014}
Gilchrist, S., Sim, J. W., and Zakrajsek, E. (2014).
Uncertainty, financial frictions, and investment dynamics.
NBER Working Paper 20038.

\bibitem{HolmstromTirole1997}
Holmstrom, B., and Tirole, J. (1997).
Financial intermediation, loanable funds, and the real sector.
\emph{Quarterly Journal of Economics}, 112(3), 663--691.

\bibitem{KiyotakiMoore1997}
Kiyotaki, N., and Moore, J. (1997).
Credit cycles. \emph{Journal of Political Economy}, 105(2), 211--248.

\bibitem{Lahura2011}
Lahura, E. (2011).
An empirical analysis of the credit-output relationship: Evidence from Peru.
BCRP Working Paper 2011-018.

\bibitem{MullerVerner2024}
Muller, K., and Verner, E. (2024).
Credit allocation and macroeconomic fluctuations.
\emph{Review of Economic Studies}, 91(6), 3645--3676.

\bibitem{PetersenRajan1994}
Petersen, M. A., and Rajan, R. G. (1994).
The benefits of lending relationships: Evidence from small business data.
\emph{Journal of Finance}, 49(1), 3--37.

\bibitem{StiglitzWeiss1981}
Stiglitz, J. E., and Weiss, A. (1981).
Credit rationing in markets with imperfect information.
\emph{American Economic Review}, 71(3), 393--410.

\bibitem{StockWatson2003}
Stock, J. H., and Watson, M. W. (2003).
Forecasting output and inflation: The role of asset prices.
\emph{Journal of Economic Literature}, 41(3), 788--829.

\end{thebibliography}

\end{document}
